\section{Conclusion} \label{sec:34-research01-conclusion}

\begin{foldable}
  While a large training set and internal access to the teacher are the prerequisites to most knowledge distillation (KD) approaches, they are rarely available due to various restrictions.
  To address these practical challenges, we propose DivBFKD, a novel method to bridge the gap for this black-box few-shot \ac{KD} problem.
  We investigate the diversity of training data, a crucial but inadequately addressed factor for \ac{KD}.
  Specifically, we propose the use of high-confidence images---ones that the teacher can classify confidently.
  We also define adaptive thresholds, the teacher-supervised criteria to negate teacher's class-wise bias during the selection of these images.
  Under the teacher's supervision, we introduce high-confidence images to the WGAN training loop on-the-fly to boost the diversity of image generation process with negligible cost.
  Through comprehensive evaluation, we have shown that our DivBFKD could use diversity to achieve state-of-the-art \ac{KD} performance among other few-shot \ac{KD} methods on multiple image datasets and architectures at different scales.
  As a future work, it would be interesting to study the robustness aspects of the distilled model relative to the original model to have the assurance~\cite{gopakumar2018algorithmic} for successful deployment in the real world.
\end{foldable}
